\documentclass[main.tex]{subfiles}
\begin{document}


\section*{Hohe Kohäsion}

\textbf{Hohe Kohäsion} bezieht sich auf das Prinzip, dass alle Elemente innerhalb einer Klasse (Methoden und Attribute) eng miteinander verbunden sind und eine gemeinsame Verantwortung oder Aufgabe haben. 

Vorteile hoher Kohäsion:
\begin{itemize}
    \item Bessere Designs: Klassen sind einfacher zu verstehen und zu verwenden.
    \item Erhöhte Wartungsfähigkeit: Änderungen an einer kohäsiven Klasse betreffen in der Regel nur die Klasse selbst.
    \item Bessere Wiederverwendbarkeit: Klassen, die klar definierte Aufgaben erfüllen, können leichter in anderen Kontexten verwendet werden.
\end{itemize}

\section*{Was bedeutet lose Kopplung?}

\textbf{Lose Kopplung} ist ein Prinzip in der Softwareentwicklung, das darauf abzielt, die Abhängigkeiten zwischen Klassen zu minimieren. 

Vorteile loser Kopplung:
\begin{itemize}
    \item Flexibilität: Klassen können unabhängig voneinander geändert, getestet und gewartet werden.
    \item Einfachere Wiederverwendbarkeit: Lose gekoppelte Klassen können leichter in anderen Projekten oder Kontexten verwendet werden.
    \item Erleichterte Wartung: Änderungen in einer Klasse erfordern keine umfangreichen Anpassungen in anderen Klassen.
\end{itemize}

\end{document}